\section{Valuations on Frames}
\label{sec:valuation}

\subsection{The Event Logic}

\begin{definition}[frame]
\label{def:frame}
A \emph{frame} is a complete lattice $\Omega$ in which finite meets
distribute over arbitrary joins,
\[
  a \sqcap \Bigl(\bigsqcup_{i \in I} b_i\Bigr)
  \;=\; \bigsqcup_{i \in I} (a \sqcap b_i)
\]
for every $a \in \Omega$ and every family $\{b_i\}_{i \in I}
\subseteq \Omega$.  Equivalently, it is a complete Heyting algebra.
By $x \Rightarrow y$ we denote the Heyting implication in a frame,
and $\neg a := a \Rightarrow \bot$ is the pseudocomplement of
$a \in \Omega$.  Also, $a^{c}$ denotes a complement of $a$, if it
exists, i.e.\ an element such that $a \sqcap a^{c} = \bot$ and
$a \sqcup a^{c} = \top$.
\end{definition}

\noindent
Frames are the algebraic semantics of
intuitionistic logic and the point-free abstraction of topological
spaces.  The same objects with morphisms reversed form the category
of \emph{locales}, so the distinction is invisible at the level of
a single algebra.  We say \emph{frame} throughout and reserve
\emph{locale} for the spatial and categorical statements (points of
a locale, classifying locales) where that word is the standard one.  The motivating example throughout is
$\Omega = \mathcal{O}(X)$, the opens of a space, with
$\sqcap, \sqcup$ as intersection and union and the pseudocomplement
$\neg a$ of an open given by the \emph{interior} of its set
complement.  In \textsf{mathlib} this is the typeclass
\texttt{Order.Frame}, and the Heyting operations come with the key
asymmetry already in place: $a \sqcap \neg a = \bot$ holds in
any Heyting algebra (\texttt{inf\_compl\_eq\_bot}), while
$a \sqcup \neg a = \top$ is available only in Boolean
algebras.  An element with $a \sqcup \neg a = \top$ is
\emph{complemented}: it satisfies its own instance of excluded
middle.

\subsection{The Valuation Structure}

The central object assigns a degree of certainty to each element of
the frame, with inclusion--exclusion as the only interaction axiom:

\begin{definition}[\texttt{Valuation}]
\label{def:valuation}
A \emph{valuation} on a frame $\Omega$ is a \emph{monotone} function
$v : \Omega \to \ENNReal$ with
$v(\bot) = 0$, $v(\top) = 1$ (\emph{normalization}), and, for all
$a, b \in \Omega$, \emph{modularity}:
\[
  v(a) + v(b) \;=\; v(a \sqcup b) + v(a \sqcap b).
\]
\end{definition}

\begin{lstlisting}
structure Valuation (Ω : Type*) [Order.Frame Ω] where
  toFun : Ω → ℝ≥0∞
  map_bot' : toFun ⊥ = 0
  map_top' : toFun ⊤ = 1
  mono'    : Monotone toFun
  modular' : ∀ a b, toFun a + toFun b
             = toFun (a ⊔ b) + toFun (a ⊓ b)
\end{lstlisting}

Values live in $\ENNReal$ (\texttt{ENNReal}, extended non-negative
reals) following \textsf{mathlib}'s measure-theoretic convention.
$\ENNReal$-valued sums are unconditionally defined for every family,
need no summability hypothesis, and commute with each other over
sigma types (\texttt{ENNReal.tsum\_sigma'}), which is exactly what
lets the countable-mixture results discussed in
Section~\ref{sec:conclusion} hold without extra side conditions.
Subtraction is truncated at $0$, but this never matters here, since it
is applied only to pairs already ordered by monotonicity, as in
\texttt{Valuation.mass} of Section~\ref{sec:representation}.  
Normalization then confines
values to $[0,1]$
(\texttt{Valuation.le\_one}).  

\subsection{Slack}

Modularity applied to the disjoint pair $(a, \neg a)$
collapses to
$v(a) + v(\neg a) = v(a \sqcup \neg a)$
(\texttt{Valuation.add\_compl\_eq\_sup}).
So, the join need not be $\top$:

\begin{theorem}[\texttt{Valuation.add\_compl\_le\_one}]
For every valuation $v$ on a frame $\Omega$ and every $a \in \Omega$,
$v(a) + v(\neg a) \le 1$.
\end{theorem}

\noindent
This is a \emph{derived} Dempster--Shafer sub-additivity inequality.
The deficit
\[
  \slack(a) \;=\; 1 - \bigl(v(a) + v(\neg a)\bigr)
  \;=\; 1 - v(a \sqcup \neg a)
\]
is the valuation of the undecided region.  It decomposes canonically.
The pair $(\neg\neg a, \neg a)$, double negation
against negation, is always disjoint. However, it is a partition only under
\emph{weak} excluded middle
($\neg a \sqcup \neg\neg a = \top$).
Applying modularity to this disjoint pair splits
the slack into two logically independent gaps:

\begin{definition}[regular element]
An element $a \in \Omega$ is \emph{regular} (or $\neg\neg$-stable) if
$a = \neg\neg a$.
\end{definition}

\begin{theorem}[\texttt{Valuation.slack\_eq\_\allowbreak dnGap\_\allowbreak add\_deMorganGap}]
For every valuation $v$ on a frame $\Omega$ and every $a \in \Omega$,
$\slack(a) = \mathrm{dnGap}(a) + \mathrm{deMorganGap}(a)$, where
$\mathrm{dnGap}(a) = v(\neg\neg a) - v(a)$ is $0$ if $a$ is
$\neg\neg$-stable (regular), and
$\mathrm{deMorganGap}(a) = 1 - v(\neg\neg a \sqcup \neg a)$
is $0$ if weak excluded middle holds at $a$.
\end{theorem}

\noindent
The single Dempster--Shafer ``ignorance'' number is thus resolved
into a regularity defect and a De~Morgan defect: the first measures
failure of double-negation stability, the second corresponds to a
named intermediate logic (weak excluded middle).  
Constructively, the contrapositive of a universal
hypothesis returns only its double-negated form, so for open
sets, this specification of probability does not fall victim
to Hempel's
raven paradox~\cite{hempel1945}: evidence of the white
sock indoors confirms at most $\neg\neg a$, not $a$ itself, and
$\mathrm{dnGap}$ is exactly the gap between confirming
$\neg\neg a$ and confirming $a$ itself.

Because $\mathrm{dnGap}$ and $\mathrm{deMorganGap}$ are themselves
defined from $v$, the apparent three unknowns $v(a)$, $v(\neg a)$, and
$v(\neg\neg a)$ are not independent: they are part of a three-constraint 
system, yielding the following result:

\begin{theorem}[\texttt{Valuation.slack\_eq\_sub\_add\_sub},
  \texttt{Valuation.le\_compl\_compl}]
\label{thm:slack-explicit}
For every valuation $v$ on a frame $\Omega$ and every $a \in \Omega$,
\[
  \slack(a) \;=\; \bigl(v(\neg\neg a) - v(a)\bigr)
  \;+\; \bigl(1 - v(\neg\neg a) - v(\neg a)\bigr),
\]
where $v(a) \le v(\neg\neg a)$.
\end{theorem}

\noindent
The inequality follows from monotonicity applied to
$a \le \neg\neg a$, a fact of every Heyting algebra, and shows the
first summand is exactly the difference, not the floor
produced by $\ENNReal$'s truncated subtraction.  So the formula above
pins $\slack(a)$ down completely from $v(a)$, $v(\neg a)$, and
$v(\neg\neg a)$ alone.  The split into a
regularity defect and a De~Morgan defect follows from monotonicity
and modularity alone, with no further axiom about ignorance required.
This contrasts with the usual Dempster--Shafer picture, where the
mass assigned to non-singleton, ambiguous sets is an extra primitive.

\begin{remark}[\texttt{Valuation.slack\_eq\_dnGap\_\allowbreak of\_isDeMorgan}]
The two gaps, $\mathrm{dnGap}$ and $\mathrm{deMorganGap}$, computed 
at a single element, are logically independent.
However, if
weak excluded middle is imposed on the frame for every $a$
($\neg\neg a \sqcup \neg a = \top$, the
condition \texttt{Order.Frame.IsDeMorgan}), then
$\mathrm{deMorganGap}$ vanishes identically and
$\slack(a) = \mathrm{dnGap}(a)$ for every valuation and every $a$.
This sits strictly between plain intuitionistic logic, where both
gaps can be positive, and full excluded middle, where both vanish
(\texttt{classical\_slack\_zero}).  
This follows from decomposition for free, and a concrete
instance appears in Section~\ref{sec:bridges}.
\end{remark}

\subsection{What Survives without Excluded Middle}

Three classical pillars survive intact, and one fails in a precisely
delimited way.  \emph{Disjoint additivity} is unconditional:
$a \sqcap b = \bot$ implies $v(a \sqcup b) = v(a) + v(b)$.
\emph{Conditioning} is native: for $v(b) \neq 0$ the posterior
$v(a \mid b) = v(a \sqcap b)/v(b)$ is again a valuation (the proof is
frame distributivity), and the product rule
$v(a \mid b) \cdot v(b) = v(a \sqcap b)$ and Bayes symmetry
(and therefore Bayesian updating)
$v(a \mid b)\,v(b) = v(b \mid a)\,v(a)$ hold
(\texttt{Valuation.condVal\_mul}, \texttt{Valuation.condVal\_symm}).  \emph{Total
probability} holds over any genuine partition, meaning
$a \sqcap a' = \bot$ and $a \sqcup a' = \top$:
$v(b) = v(a \sqcap b) + v(a' \sqcap b)$.  What fails is total
probability over $\{a, \neg a\}$.  That family does not tile
$\top$, and the conditional masses fall short of $v(b)$ by a
conditional slack (\texttt{Valuation.cond\_add\_compl\_le}).

\paragraph{The classical fragment.}
We characterize the elements of $\Omega$ on which the classical
complement rule holds.  The $\neg\neg$-stable (\emph{regular})
elements, which form a Boolean algebra (by Glivenko's theorem), are a
natural candidate, but they do not suffice. Modularity constrains
the frame join $\sqcup$, whereas the Boolean join of regular
elements is the regularized $\neg\neg(a \sqcup b)$, for which
only the inequality
$v(a) + v(b) \le v(\neg\neg(a \sqcup b)) + v(a \sqcap b)$ survives
(i.e. the 2-monotonicity of Section~\ref{sec:ds}).  A regular but uncomplemented
element, which is an open ray $(-\infty,0)$ in the opens of $\mathbb{R}$,
under an atomic measure at $0$, still carries slack.  The correct
classical fragment is the \emph{complemented} elements, where
$v(a) + v(a^{c}) = 1$ holds exactly
(\texttt{Valuation.add\_compl\_eq\_one\_of\_\allowbreak complemented}).  On a connected
locale that fragment can be as small as $\{\bot, \top\}$.  In the
Boolean limit, every element is complemented, the slack vanishes
identically, and Kolmogorov is recovered
(\texttt{classical\_additivity}, \texttt{classical\_slack\_zero}).

\paragraph{Mixtures.}
Valuations are closed under finite convex combination
(\texttt{Valuation.mix}), with every axiom inherited termwise,
because modularity is linear.  This innocuous fact becomes a
characterization in Section~\ref{sec:representation}.
